\makelecturetitle{Множества и мощность: введение в тему}

\tableofcontents

\section{Введение}

\term{Множество} — это совокупность объектов, рассматриваемых как единое целое.
Мощность множества $|A|$ показывает, сколько в нём элементов.

\begin{definition}
Множества $A$ и $B$ называются \term{равномощными}, если между ними существует
биекция: каждому элементу $A$ соответствует ровно один элемент $B$ и наоборот.
\end{definition}

\begin{theorem}
Для любых конечных множеств $A$ и $B$ справедлива формула включений-исключений:
\[ |A \cup B| = |A| + |B| - |A \cap B|. \]
\end{theorem}

\begin{proof}
Каждый элемент $A \cap B$ посчитан дважды в сумме $|A| + |B|$, поэтому его нужно
вычесть один раз. Остальные элементы входят ровно в одно из множеств.
\qed
\end{proof}

\section{Примеры и задачи}

\begin{example}
Мощность множества цифр $\{0, 1, \ldots, 9\}$ равна 10, а множества гласных
русского алфавита — 10.
\end{example}

\begin{important}
Не путайте мощность множества с мощностью системы множеств: $|A|$ и $|\mathcal{P}(A)|$
считают разные вещи. Для конечного $A$ верно $|\mathcal{P}(A)| = 2^{|A|}$.
\end{important}

\begin{table}[h]
\centering
\begin{tabularx}{0.8\linewidth}{lX}
\toprule
Операция & Смысл \\
\midrule
$A \cup B$ & объединение: элементы хотя бы одного множества \\
$A \cap B$ & пересечение: общие элементы \\
$A \setminus B$ & элементы $A$, которых нет в $B$ \\
\bottomrule
\end{tabularx}
\caption{Основные операции над множествами}
\end{table}

\begin{enumerate}
  \item Найдите мощность множества $\{x \in \mathbb{N} \mid x < 10\}$.
  \item Докажите формулу включений-исключений для трёх множеств.
  \item Приведите практический пример из вашей предметной области.
\end{enumerate}

\section{Выводы}

\begin{keyidea}
Мы освоили базовые определения и научились применять их на практике.
На следующем занятии перейдём к счётным и несчётным множествам.
\end{keyidea}
